\chapter{Crisis Management}\label{ch:fm:crisis-management}

On 31 October 2011 liquidation proceedings began for MF Global, a large futures commission merchant, and its customers learned that money which
should have been held apart for them was missing. The regulator later found that during the last week of October the firm had used customer
segregated funds to support its own and its affiliates' operations; a court ordered \$1.212 billion of restitution in 2013, and customers were
eventually repaid in full. In the first days, firms that cleared through it had to find out which of their positions still existed, where their
margin was and who could move it. A crisis gives a firm hours, and what it can do in those hours is decided long before.

\section{The first hour of a technology incident}

\begin{definition}[Crisis management plan, crisis committee]\label{def:fm:crisis-management:plan}
A \emph{crisis management plan}\index{crisis management plan} sets out, for defined kinds of crisis (a trading system out of control, a liquidity
squeeze, a counterparty failure, a market closure), who decides, what they may do, whom they tell and in what order, with the runbooks and contact
lists to do it. The \emph{crisis committee}\index{crisis committee} is the small group the plan names to take those decisions, with the authority to
stop trading, cut positions, move collateral and draw funding without the usual approvals.
\end{definition}

The first hour of a technology incident is Book 13's and Book 15's: the kill switch and the mass cancel stop the damage, the incident commander
runs the response from the runbook, failover moves trading to the disaster-recovery site (Book 14) if the primary is lost. What the \omterm{def:fm:crisis-management:plan}{crisis committee}
adds is the business decision around the technical one: whether to resume, which positions to flatten, what to tell counterparties, clients and the
regulator, and how much cash the incident has already cost. Book 14's business continuity plan and operational resilience rules set the recovery
targets; the crisis plan decides who acts when those targets are about to be missed.

\section{Liquidity squeezes: margin calls you cannot meet}

A liquidity crisis is chapter 14's stress test happening in hours instead of days: requirements rise at every broker and clearing house at once,
funding sources that looked reliable are withdrawn, and assets that could be sold on a normal day take longer or cost more. The \omterm{def:fm:treasury-and-funding:horizon}{survival horizon} is
measured in hours, and every action has a delay.

\begin{definition}[Liquidity drill]\label{def:fm:crisis-management:drill}
A \emph{liquidity drill}\index{liquidity drill} rehearses a liquidity crisis on the firm's real accounts, terms and funding sources: a scripted
scenario is run hour by hour, the \omterm{def:fm:crisis-management:plan}{crisis committee} chooses actions under their real delays and costs, and the drill records the \omterm{def:fm:treasury-and-funding:horizon}{survival horizon}
and the decisions, to be compared with the plan.
\end{definition}

\section{When a counterparty fails}

\begin{definition}[Trapped assets]\label{def:fm:crisis-management:trapped}
\emph{Trapped assets}\index{trapped assets} are a firm's cash or securities held by a counterparty that has failed or been placed in insolvency
proceedings, which the firm cannot use until an administrator or trustee returns them, if ever: collateral posted, excess margin, securities
rehypothecated (Book 1), or balances not segregated as the rules required.
\end{definition}

When a prime broker or clearing broker fails, three things happen at once. The collateral it held is trapped, wholly or in part, for months or
years (Book 1 recounted the administration of Lehman Brothers' European broker-dealer). The positions it held must be moved or re-established
elsewhere, and the new broker requires margin for them immediately. And the firm's own counterparties, seeing the news, may raise their terms.
Customer asset segregation (Book 3) is the protection; MF Global showed that its value depends on the broker obeying it.

\begin{dated}{2026-09}{The public record of the chapter's examples}\label{dat:fm:crisis-management:record}
\textbf{MF Global}: liquidation proceedings began on 31 October 2011; the CFTC's complaint charged that in the last week of October 2011 the firm
used customer segregated funds to support its own and its affiliates' operations; a consent order of 8 November 2013 required MF Global Inc. to pay
\$1.212 billion of restitution to customers and a \$100 million penalty, and the trustee later made final distributions satisfying full
restitution. \textbf{Gilt market, 2022}: on 28 September 2022 the Bank of England announced temporary purchases of long-dated UK government bonds
to restore orderly market conditions, after its Financial Policy Committee noted the risks to financial stability from dysfunction in that market;
the auctions ran until 14 October.
\end{dated}

\section{Tutorial: seventy-two hours}\label{tut:fm:crisis-management:tut}

\begin{tutorial}
\textbf{Goal.} Run a 72-hour \omterm{def:fm:crisis-management:drill}{liquidity drill} on the firm of chapter 14's kind: a volatility shock, a prime broker that fails, collateral trapped; find
the \omterm{def:fm:treasury-and-funding:horizon}{survival horizon} and the cheapest set of actions that extends it past 72 hours. \textbf{End state:} the hourly cash chart with and without the
actions (\cref{fig:fm:crisis-management:cash}) and the plan's steps (\cref{fig:fm:crisis-management:steps}).
\begin{tutsteps}
\item \textbf{The accounts.} Prime broker A holds \$44 million against a requirement of \$40 million; prime broker B \$33 million against \$30 million;
the clearing broker \$22 million against \$20 million. \omterm{def:fm:treasury-and-funding:cash}{Unencumbered cash} is \$30 million.
\item \textbf{The shock.} \omterm{def:fm:treasury-and-funding:house}{House margin} rises by 60\% over the first day (\texttt{firm.treasury.house\_path}, no lock-up); the clearing house raises
initial margin by 30\% at the first daily cycle and 50\% at the second.
\item \textbf{The failure.} Prime broker B fails at hour 6: 60\% of the \$33 million it holds is trapped, the rest comes back at hour 48, and its
positions move to prime broker A.
\item \textbf{The actions.} Sell money-market funds (\$10 million after 4 hours, cost \$0.02 million); repo the bond portfolio (\$10 million after 8
hours, \$0.1 million); draw the committed line (\$20 million after 2 hours, \$0.3 million, withdrawn at hour 12); cut positions by 10\%, 20\% or 30\%
(after 3 hours, at impact costs of \$1, 3 and 6 million).
\item \textbf{The plan.} \texttt{firm.crisisdrill.greedy} adds the action with the lowest cost per hour gained until the horizon reaches 72 hours.
\end{tutsteps}
\end{tutorial}

\omcode{code/firm/crisisdrill/firm_crisisdrill.py}{52}{90}{The hourly liquidity engine: calls paid from cash as requirements rise, the failed broker's
positions moved and its collateral trapped, and actions completing after their delays.}\label{lst:fm:crisis-management:run}

\begin{omfigure}
\begin{tikzpicture}
\begin{axis}[width=10.5cm,height=5.4cm,xlabel={\small hour},ylabel={\small unencumbered cash (\$ million)},xmin=0,xmax=71,ymin=-50,ymax=75,
  tick label style={font=\footnotesize},grid=major,grid style={gray!20},table/col sep=comma,
  legend cell align=left,legend style={font=\scriptsize,at={(0.5,-0.3)},anchor=north,/tikz/every even column/.append style={column sep=8pt}},legend columns=3]
\addplot[omDef,very thick,const plot] table[x=hour,y=baseline]{figdata/desk/27-crisis-management/cash.csv};
\addplot[omThm,thick,const plot] table[x=hour,y=plan]{figdata/desk/27-crisis-management/cash.csv};
\addplot[gray,dashed,const plot] table[x=hour,y=nofail]{figdata/desk/27-crisis-management/cash.csv};
\draw[black,thin] (axis cs:0,0) -- (axis cs:71,0);
\draw[gray,dotted] (axis cs:6,-50) -- (axis cs:6,75) node[pos=0.05,right,font=\scriptsize] {broker fails};
\legend{no action,cheapest plan,no failure}
\end{axis}
\end{tikzpicture}

\omcaption{Seventy-two hours of \omterm{def:fm:treasury-and-funding:cash}{unencumbered cash} under the drill: with no action the firm runs out of cash in hour 6, when the failed broker's positions
must be margined again at the surviving one; with the cheapest plan it keeps at least \$11.8 million. Without the failure the shock alone exhausts the
cash in hour 23. Data: \texttt{fm\_crisis}.}\label{fig:fm:crisis-management:cash}
\end{omfigure}

Without action the firm is out of cash in hour 6. The failure does the damage: \$19.8 million of collateral is trapped at the failed broker, and its
positions need margin at the surviving one at once; the shock alone would have exhausted the cash in hour 23. Over the 72 hours the calls total \$78.8
million. The greedy plan sells the money-market funds first (the horizon moves to hour 9 for \$20\,000), then repos the bonds (hour 15), then draws the
line before it is withdrawn (hour 24), and only then cuts 10\% of positions, which carries the firm past 72 hours: \$1.42 million in all, most of it
the cut's market impact. The three cash sources together reach only hour 24; without the cut, the second daily clearing cycle would still exhaust the
cash.

\begin{omfigure}
\begin{tikzpicture}
\begin{axis}[width=10.5cm,height=5cm,xbar,bar width=9pt,xmin=0,xmax=80,xlabel={\small survival horizon (hours)},ytick={0,1,2,3,4},
  yticklabels={no action,{+ sell money-market funds},{+ repo the bonds},{+ draw the line},{+ cut 10\% of positions}},y dir=reverse,
  ymin=-0.6,ymax=4.6,tick label style={font=\footnotesize},grid=major,grid style={gray!20},table/col sep=comma,nodes near coords,
  nodes near coords style={font=\scriptsize}]
\addplot[fill=omDef!55,draw=omDef] table[y=pos,x=horizon]{figdata/desk/27-crisis-management/steps.csv};
\draw[gray,dashed] (axis cs:72,-0.6) -- (axis cs:72,4.6);
\end{axis}
\end{tikzpicture}

\omcaption{The \omterm{def:fm:treasury-and-funding:horizon}{survival horizon} as the greedy plan adds actions, cheapest per hour gained first; cumulative costs \$0, 0.02, 0.12, 0.42 and 1.42 million.
Data: \texttt{fm\_crisis.steps}.}\label{fig:fm:crisis-management:steps}
\end{omfigure}

\begin{proposition}[More resources never shorten the horizon]\label{prop:fm:crisis-management:monotone}
In the drill's model, adding an action that brings cash or cuts requirements, without changing any other action, never makes the \omterm{def:fm:treasury-and-funding:horizon}{survival horizon}
shorter.
\end{proposition}

\begin{proof}
Cash at each hour is the starting cash plus completed actions' cash, plus returns of excess, less calls. A cash action adds a non-negative amount
from its completion onwards and changes no call. A cut lowers every later requirement, so every later call is no larger (the prime brokers retain
excess, and the clearing account returns it, which only adds cash). Cash is therefore at least as high at every hour, and the first negative hour is no
earlier.
\end{proof}

\begin{remark}[Greedy is not optimal]\label{rem:fm:crisis-management:greedy}
Choosing by cost per hour gained is a heuristic. Hours are not additive: an action with a long delay gains nothing alone and a great deal after
others have bought the time it needs, and a line that is withdrawn at hour 12 must be drawn early or never. For a plan with a handful of actions, an
exhaustive search is cheap; the drill's value is in the committee practising the choice, not in the algorithm.
\end{remark}

\section{Drills and playbooks}

\begin{method}[Running a liquidity drill]\label{met:fm:crisis-management:drill}
\begin{enumerate}
\item Write the scenario from the firm's own accounts and terms: which broker fails, what is trapped, what is withdrawn, how margins rise, hour by
hour.
\item Give the \omterm{def:fm:crisis-management:plan}{crisis committee} the real contact lists, the real authorities and the real delays of each action, and run it without warning.
\item Record every decision with its time; compute the \omterm{def:fm:treasury-and-funding:horizon}{survival horizon} achieved against the plan's.
\item Fix what the drill exposed: a line that needed two signatures nobody could reach, a transfer that missed a cut-off, a runbook that assumed a
system that was down.
\item Repeat at least yearly, and after any change of broker, funding line or strategy.
\end{enumerate}
\end{method}

The Financial Stability Board's 2024 recommendations on liquidity preparedness for margin and collateral calls, summarised in chapter 14, ask
non-bank firms for exactly this: stress tests of margin calls under extreme but plausible scenarios and plans to meet them. A drill is where the
plan meets the clock.

\section{Build: the crisis drill}\label{bld:fm:crisis-management:crisisdrill}

\begin{build}
\bfield{Purpose} An hourly \omterm{def:fm:crisis-management:drill}{liquidity drill}: scenarios with shocks, a failing counterparty, \omterm{def:fm:crisis-management:trapped}{trapped assets} and withdrawn lines; the \omterm{def:fm:treasury-and-funding:horizon}{survival horizon};
and an action plan chosen by cost per hour gained.
\bfield{Interface} \texttt{firm.crisisdrill}: \texttt{Account}, \texttt{Scenario}, \texttt{Action}, \texttt{run}, \texttt{greedy}; the house-margin path
from \texttt{firm.treasury}.
\bfield{Rules} Calls are paid the hour they arise; a failed account stops and returns its untrapped collateral at the stated hour; an action is
available only up to its deadline.
\bfield{Acceptance tests} \texttt{code/firm/crisisdrill/tests/}: calls and horizon on hand numbers; failure, trapping and transfer; the greedy choice
and a withdrawn line.
\bfield{Stretch} Exhaustive search over action sets; uncertain delays; the clearing house's own margin model from \texttt{firm.initmargin}; several
failures.
\end{build}

\begin{omsources}
\item CFTC press releases 6776-13 (2013), 6904-14 (2014) and 7508-17 (2017) on MF Global.
\item Bank of England, ``Bank of England announces gilt market operation'', 28 September 2022.
\end{omsources}

\section{Exercises}

\begin{exercise}[$\star$]\label{exo:fm:crisis-management:1}
How much collateral is trapped at the failed broker, and how much comes back at hour 48?
\end{exercise}

\begin{exercise}[$\star$]\label{exo:fm:crisis-management:2}
Define \omterm{def:fm:crisis-management:trapped}{trapped assets} and give three ways a firm's assets can be trapped at a failed broker.
\end{exercise}

\begin{exercise}[$\star$]\label{exo:fm:crisis-management:3}
In which hour does the firm run out of cash without the failure, and why so much later?
\end{exercise}

\begin{exercise}[$\star\star$]\label{exo:fm:crisis-management:4}
Why does the greedy plan draw the committed line only third, and what would happen if the line needed two hours of approvals from a person who is
unreachable until hour 12?
\end{exercise}

\begin{exercise}[$\star\star$]\label{exo:fm:crisis-management:5}
Prove \cref{prop:fm:crisis-management:monotone}. Would it hold if cutting positions triggered losses paid in cash?
\end{exercise}

\begin{exercise}[$\star\star$]\label{exo:fm:crisis-management:6}
What did MF Global's customers need to know in the first days, and which records would have told them?
\end{exercise}

\begin{exercise}[$\star\star\star$]\label{exo:fm:crisis-management:7}
\emph{Coding.} Run the drill with all three cash actions and no cut. What is the horizon, and what is the lowest cash?
\end{exercise}

\begin{exercise}[$\star\star\star$]\label{exo:fm:crisis-management:8}
\emph{Find the flaw.} ``We have \$30 million of cash and \$40 million of lines: we can survive any margin call.''
\end{exercise}

\section{Problem: Seventy-Two Hours}

\begin{problem}[{Weekend problem --- seventy-two hours}]\label{pb:fm:crisis-management:1}
A \omterm{def:fm:the-risk-management-function:cro}{chief risk officer} must run the firm's first \omterm{def:fm:crisis-management:drill}{liquidity drill} and report what it found to the board.

\medskip\noindent\textbf{Part I --- The plan.}
\begin{enumerate}
\item Define a \omterm{def:fm:crisis-management:plan}{crisis management plan} and a \omterm{def:fm:crisis-management:plan}{crisis committee}.
\item What does the committee add to a technology incident's response?
\item Define a \omterm{def:fm:crisis-management:drill}{liquidity drill}.
\item Define \omterm{def:fm:crisis-management:trapped}{trapped assets}.
\end{enumerate}

\medskip\noindent\textbf{Part II --- The record.}
\begin{enumerate}[resume]
\item Summarise the public record on MF Global.
\item What did the Bank of England do in September 2022, and why?
\item What protects customer assets at a broker, and what did MF Global show about it?
\item What do the 2024 international recommendations ask of non-bank firms?
\end{enumerate}

\medskip\noindent\textbf{Part III --- The drill.}
\begin{enumerate}[resume]
\item Describe the scenario hour by hour.
\item Give the \omterm{def:fm:treasury-and-funding:horizon}{survival horizon} with no action, and without the failure.
\item Give the greedy plan's steps, horizons and costs.
\item State and prove \cref{prop:fm:crisis-management:monotone}.
\item Why is the greedy plan not necessarily the cheapest?
\end{enumerate}

\medskip\noindent\textbf{Part IV --- The report.}
\begin{enumerate}[resume]
\item Which weaknesses would a real drill likely expose that the model does not?
\item What would you change in the firm's broker arrangements?
\item How large should the \omterm{def:fm:treasury-and-funding:cash}{liquidity buffer} be?
\item How often should the drill run, and who should not know in advance?
\item What would you report to the board?
\item State the \emph{named result}: the \omterm{def:fm:treasury-and-funding:horizon}{survival horizon} under the drill, and the cost of the cheapest action set that extends it to 72 hours.
\item In two sentences, write the recommendation.
\end{enumerate}
\end{problem}

\section{Interview questions}

\begin{interviewq}[$\star$ \iqroles{developer}]\label{iq:fm:crisis-management:1}
A strategy is sending orders it should not. What do you do in the first five minutes?
\end{interviewq}

\begin{interviewq}[$\star$ \iqroles{risk}]\label{iq:fm:crisis-management:2}
Your prime broker has just announced it is in administration. What do you check first?
\end{interviewq}

\begin{interviewq}[$\star\star$ \iqroles{risk, trader}]\label{iq:fm:crisis-management:3}
Margin calls exceed your cash by \$10 million tomorrow morning. List your options in order.
\end{interviewq}

\begin{interviewq}[$\star\star$ \iqroles{developer, risk}]\label{iq:fm:crisis-management:4}
How would you design a \omterm{def:fm:crisis-management:drill}{liquidity drill} that tests people as well as models?
\end{interviewq}

\begin{interviewq}[$\star\star$ \iqroles{trader}]\label{iq:fm:crisis-management:5}
Why can cutting positions in a crisis cost more than it saves?
\end{interviewq}

\begin{interviewq}[$\star\star\star$ \iqroles{researcher, risk}]\label{iq:fm:crisis-management:6}
Formulate the choice of crisis actions as an optimisation problem, and say why a greedy rule can fail.
\end{interviewq}
